% GENERATED FILE. Edit canonical lessons, then run npm run generate:books.
% canonical-source-hash: fnv1a64:5a9873aa06fd3d70
% canonical-lessons: ML-C352-patrasammelanam, ML-C352-prastavana, ML-C352-satyavastha, ML-C352-peruppikkuka, ML-C352-mappu-parayuka

\chapter{A Press Conference}
\label{ch:ml-a-press-conference}
\hlchaptermodality{352}

\begin{chapteropening}
\textbf{By the end of this chapter:} \emph{I can say a press conference, a statement, the real facts, to exaggerate and to say sorry.}

Complete the last lesson of chapter 352: I can say a press conference, a statement, the real facts, to exaggerate and to say sorry.
\end{chapteropening}

\section[\texorpdfstring{\ml{പത്രസമ്മേളനം} (patrasammēḷanaṁ)}{പത്രസമ്മേളനം (patrasammēḷanaṁ)}]{\ml{പത്രസമ്മേളനം} (patrasammēḷanaṁ) — a press conference}

\label{lesson:ML-C352-patrasammelanam}



\begin{quote}
Before the new one: say the Malayalam for to depict, then the Malayalam for to broadcast.
\end{quote}

\subsection*{You'll want to know: \ml{പത്രസമ്മേളനം}}
\textbf{\ml{പത്രസമ്മേളനം}} — \emph{patrasammēḷanaṁ} — \textquotedblleft{}a press conference\textquotedblright{}.

\begin{quote}
\ml{മാനേജർ} \ml{ഇന്ന്} \ml{ഒരു} \ml{പത്രസമ്മേളനം} \ml{നടത്തി}.
\end{quote}

\begin{quote}
\emph{mānējar innŭ oru patrasammēḷanaṁ naṭatti.}
\end{quote}

\begin{quote}
The manager held a press conference today.
\end{quote}

A press conference is held: \ml{പത്രസമ്മേളനം} \ml{നടത്തുക}.

\begin{cousinweb}
\textbf{\ml{പത്രസമ്മേളനം}} (\emph{patrasammēḷanaṁ}) is built from \ml{പത്രം} (\emph{patraṁ}), \textquotedblleft{}a newspaper\textquotedblright{}, and Sanskrit \emph{sammelana}, \textquotedblleft{}a meeting together, a union\textquotedblright{} (Monier-Williams): a meeting called for the papers.
\end{cousinweb}

\begin{grammarlens}[title={What you've built}]
One more word for telling what happened.
\end{grammarlens}

\subsection*{Your turn}
\begin{itemize}
\raggedright
  \item \emph{Say it:} \emph{patrasammēḷanaṁ}
  \item \emph{Say it:} \emph{patrasammēḷanaṁ}, once more
  \item \emph{Say it:} \emph{mānējar innŭ oru patrasammēḷanaṁ naṭatti.} — The manager held a press conference today.
  \item \emph{Recall:} say the Malayalam for to depict, then the Malayalam for to broadcast, then say \emph{patrasammēḷanaṁ} again
\end{itemize}

\begin{tcolorbox}[breakable,colback=teal!4,colframe=teal!35!black,title={Before you move on}]
What does \ml{പത്രസമ്മേളനം} mean? (\textquotedblleft{}A press conference\textquotedblright{}.) Say it once more.
\end{tcolorbox}
\section[\texorpdfstring{\ml{പ്രസ്താവന} (prastāvana)}{പ്രസ്താവന (prastāvana)}]{\ml{പ്രസ്താവന} (prastāvana) — a statement}

\label{lesson:ML-C352-prastavana}



\begin{quote}
Before the new one: say the Malayalam for to broadcast, then the Malayalam for a press conference.
\end{quote}

\subsection*{You'll want to know: \ml{പ്രസ്താവന}}
\textbf{\ml{പ്രസ്താവന}} — \emph{prastāvana} — \textquotedblleft{}a statement\textquotedblright{}.

\begin{quote}
\ml{പോലീസ്} \ml{അപകടത്തെക്കുറിച്ച്} \ml{ഒരു} \ml{പ്രസ്താവന} \ml{ഇറക്കി}.
\end{quote}

\begin{quote}
\emph{pōlīsŭ apakaṭattekkuṟiccŭ oru prastāvana iṟakki.}
\end{quote}

\begin{quote}
The police issued a statement about the accident.
\end{quote}

The formal word for what someone sets out in public, usually in writing.

\begin{cousinweb}
\textbf{\ml{പ്രസ്താവന}} (\emph{prastāvana}) is borrowed from Sanskrit \emph{prastāvanā}, \textquotedblleft{}a sounding forth; an introduction, a preface\textquotedblright{} (Monier-Williams). In Malayalam it came to mean what someone sets forth in public.
\end{cousinweb}

\begin{grammarlens}[title={What you've built}]
One more word for telling what happened.
\end{grammarlens}

\subsection*{Your turn}
\begin{itemize}
\raggedright
  \item \emph{Say it:} \emph{prastāvana}
  \item \emph{Say it:} \emph{prastāvana}, once more
  \item \emph{Say it:} \emph{pōlīsŭ apakaṭattekkuṟiccŭ oru prastāvana iṟakki.} — The police issued a statement about the accident.
  \item \emph{Recall:} say the Malayalam for to broadcast, then the Malayalam for a press conference, then say \emph{prastāvana} again
\end{itemize}

\begin{tcolorbox}[breakable,colback=teal!4,colframe=teal!35!black,title={Before you move on}]
What does \ml{പ്രസ്താവന} mean? (\textquotedblleft{}A statement\textquotedblright{}.) Say it once more.
\end{tcolorbox}
\section[\texorpdfstring{\ml{സത്യാവസ്ഥ} (satyāvastha)}{സത്യാവസ്ഥ (satyāvastha)}]{\ml{സത്യാവസ്ഥ} (satyāvastha) — the real facts; the actual truth}

\label{lesson:ML-C352-satyavastha}



\begin{quote}
Before the new one: say the Malayalam for a press conference, then the Malayalam for a statement.
\end{quote}

\subsection*{You'll want to know: \ml{സത്യാവസ്ഥ}}
\textbf{\ml{സത്യാവസ്ഥ}} — \emph{satyāvastha} — \textquotedblleft{}the real facts; the actual truth\textquotedblright{}.

\begin{quote}
\ml{ലേഖകൻ} \ml{സംഭവത്തിന്റെ} \ml{സത്യാവസ്ഥ} \ml{കണ്ടെത്തി}.
\end{quote}

\begin{quote}
\emph{lēkhakan sambhavattinṟe satyāvastha kaṇṭetti.}
\end{quote}

\begin{quote}
The correspondent found out the real facts of the incident.
\end{quote}

Literally the state of the truth: \ml{സത്യം}, truth, + \ml{അവസ്ഥ}, a state.

\begin{cousinweb}
\textbf{\ml{സത്യാവസ്ഥ}} (\emph{satyāvastha}) is built from two Sanskrit words: \emph{satya}, \textquotedblleft{}true; the truth\textquotedblright{} (CDIAL 13112), and \emph{avasthā}, \textquotedblleft{}a state, a condition\textquotedblright{} (CDIAL 872). It names how things really stand, as against what people say.
\end{cousinweb}

\begin{grammarlens}[title={What you've built}]
One more word for telling what happened.
\end{grammarlens}

\subsection*{Your turn}
\begin{itemize}
\raggedright
  \item \emph{Say it:} \emph{satyāvastha}
  \item \emph{Say it:} \emph{satyāvastha}, once more
  \item \emph{Say it:} \emph{lēkhakan sambhavattinṟe satyāvastha kaṇṭetti.} — The correspondent found out the real facts of the incident.
  \item \emph{Recall:} say the Malayalam for a press conference, then the Malayalam for a statement, then say \emph{satyāvastha} again
\end{itemize}

\begin{tcolorbox}[breakable,colback=teal!4,colframe=teal!35!black,title={Before you move on}]
What does \ml{സത്യാവസ്ഥ} mean? (\textquotedblleft{}The real facts; the actual truth\textquotedblright{}.) Say it once more.
\end{tcolorbox}
\section[\texorpdfstring{\ml{പെരുപ്പിക്കുക} (peruppikkuka)}{പെരുപ്പിക്കുക (peruppikkuka)}]{\ml{പെരുപ്പിക്കുക} (peruppikkuka) — to exaggerate}

\label{lesson:ML-C352-peruppikkuka}



\begin{quote}
Before the new one: say the Malayalam for a statement, then the Malayalam for the real facts.
\end{quote}

\subsection*{You'll want to know: \ml{പെരുപ്പിക്കുക}}
\textbf{\ml{പെരുപ്പിക്കുക}} — \emph{peruppikkuka} — \textquotedblleft{}to exaggerate\textquotedblright{}.

\begin{quote}
\ml{അവൻ} \ml{എപ്പോഴും} \ml{കഥകൾ} \ml{പെരുപ്പിച്ച്} \ml{പറയും}.
\end{quote}

\begin{quote}
\emph{avan eppōḻuṁ kathakaḷ peruppiccŭ paṟayuṁ.}
\end{quote}

\begin{quote}
He always makes his stories bigger than they were.
\end{quote}

Past: \ml{പെരുപ്പിച്ചു}, \emph{peruppiccu}.

Often with \ml{പറയുക}: \ml{പെരുപ്പിച്ച്} \ml{പറയുക}, to tell with exaggeration.

\begin{cousinweb}
\textbf{\ml{പെരുപ്പിക്കുക}} (\emph{peruppikkuka}) is native Dravidian, a causative built on \ml{പെരു} (\emph{peru}), \textquotedblleft{}great, large\textquotedblright{} (DEDR 4411): to exaggerate is to make a thing bigger than it was.
\end{cousinweb}

\begin{grammarlens}[title={What you've built}]
One more word for telling what happened.
\end{grammarlens}

\subsection*{Your turn}
\begin{itemize}
\raggedright
  \item \emph{Say it:} \emph{peruppikkuka}
  \item \emph{Say it:} \emph{peruppikkuka}, once more
  \item \emph{Say it:} \emph{avan eppōḻuṁ kathakaḷ peruppiccŭ paṟayuṁ.} — He always makes his stories bigger than they were.
  \item \emph{Recall:} say the Malayalam for a statement, then the Malayalam for the real facts, then say \emph{peruppikkuka} again
\end{itemize}

\begin{tcolorbox}[breakable,colback=teal!4,colframe=teal!35!black,title={Before you move on}]
What does \ml{പെരുപ്പിക്കുക} mean? (\textquotedblleft{}To exaggerate\textquotedblright{}.) Say it once more.
\end{tcolorbox}
\section[\texorpdfstring{\ml{മാപ്പ്} \ml{പറയുക} (māppŭ paṟayuka)}{മാപ്പ് പറയുക (māppŭ paṟayuka)}]{\ml{മാപ്പ്} \ml{പറയുക} (māppŭ paṟayuka) — to say sorry; to apologise}

\label{lesson:ML-C352-mappu-parayuka}



\begin{quote}
Before the new one: say the Malayalam for the real facts, then the Malayalam for to exaggerate.
\end{quote}

\subsection*{You'll want to know: \ml{മാപ്പ്} \ml{പറയുക}}
\textbf{\ml{മാപ്പ്} \ml{പറയുക}} — \emph{māppŭ paṟayuka} — \textquotedblleft{}to say sorry; to apologise\textquotedblright{}.

\begin{quote}
\ml{തെറ്റിന്} \ml{അവൻ} \ml{എല്ലാവരോടും} \ml{മാപ്പ്} \ml{പറഞ്ഞു}.
\end{quote}

\begin{quote}
\emph{teṟṟinŭ avan ellāvarōṭuṁ māppŭ paṟaññu.}
\end{quote}

\begin{quote}
He apologised to everyone for the mistake.
\end{quote}

Past: \ml{മാപ്പ്} \ml{പറഞ്ഞു}, \emph{māppŭ paṟaññu}.

The person takes -\ml{ഓട്}: \ml{അവളോട്} \ml{മാപ്പ്} \ml{പറഞ്ഞു}, apologised to her. \ml{മാപ്പ്} alone is pardon.

\begin{cousinweb}
\textbf{\ml{മാപ്പ്} \ml{പറയുക}} (\emph{māppŭ paṟayuka}) joins \ml{മാപ്പ്} (\emph{māppŭ}), \textquotedblleft{}pardon\textquotedblright{}, and the native verb \ml{പറയുക}, \textquotedblleft{}to say\textquotedblright{} (DEDR 4031): to apologise is to say pardon. \ml{മാപ്പ്} is usually traced to Hindustani \emph{māf}, \textquotedblleft{}forgiven\textquotedblright{}, a word of Arabic origin.
\end{cousinweb}

\begin{grammarlens}[title={What you've built}]
One more word for telling what happened.
\end{grammarlens}

\subsection*{Your turn}
\begin{itemize}
\raggedright
  \item \emph{Say it:} \emph{māppŭ paṟayuka}
  \item \emph{Say it:} \emph{māppŭ paṟayuka}, once more
  \item \emph{Say it:} \emph{teṟṟinŭ avan ellāvarōṭuṁ māppŭ paṟaññu.} — He apologised to everyone for the mistake.
  \item \emph{Recall:} say the Malayalam for the real facts, then the Malayalam for to exaggerate, then say \emph{māppŭ paṟayuka} again
\end{itemize}

\begin{tcolorbox}[breakable,colback=teal!4,colframe=teal!35!black,title={Before you move on}]
What does \ml{മാപ്പ്} \ml{പറയുക} mean? (\textquotedblleft{}To say sorry; to apologise\textquotedblright{}.) Say it once more.
\end{tcolorbox}
